% Appendix: Frobenius algebra --- deferred proofs
% Completes the well-definedness argument deferred from Proposition 4.7
% (Frob-to-TQFT), proves mapping-class invariance, and computes
% handle-operator spectra for finite-group examples.

\label{app:frobenius}

\section{From Frobenius algebras to TQFTs: well-definedness}\label{appsec:frob-to-tqft}

The forward direction of the classification theorem (Proposition~\ref{prop:TQFT-to-Frob}) was proved in full in the main text: every 2D TQFT produces a commutative Frobenius algebra. The converse direction (Proposition~\ref{prop:Frob-to-TQFT}) constructed a candidate functor $Z_V$ on the generators of $\mathrm{Cob}_2$ and verified that the local relations of Proposition~\ref{prop:cob2-relations} are satisfied. What remains is to prove that $Z_V$ is independent of the choice of Morse decomposition used to present a cobordism.

\subsection{Cerf theory and the space of Morse functions}

Let $M : \Sigma_{\mathrm{in}} \to \Sigma_{\mathrm{out}}$ be a compact oriented 2-cobordism. A \emph{Morse presentation} of $M$ is a choice of Morse function $f : M \to [0,1]$, constant on the boundary components, with distinct critical values. Different Morse functions yield different decompositions of $M$ into elementary generators.

\begin{theorem}[Cerf \cite{Cerf1970}]\label{thm:cerf}
Let $f_0, f_1 : M \to [0,1]$ be two Morse functions on a compact cobordism $M$, each with distinct critical values and constant on the boundary. Then there exists a generic one-parameter family $\{f_t\}_{t \in [0,1]}$ connecting $f_0$ to $f_1$ such that for all but finitely many values of $t$, the function $f_t$ is Morse with distinct critical values. At the finitely many exceptional values, exactly one of the following transitions occurs:
\begin{enumerate}
\item A \emph{critical-point interchange}: two critical values cross without changing the topology of the cobordism.
\item A \emph{birth--death transition}: a cancelling pair of adjacent critical points of index $k$ and $k+1$ appears or disappears.
\end{enumerate}
\end{theorem}

\noindent The theorem reduces the well-definedness problem to a finite list of local verifications: one must show that $Z_V$ is invariant under each type of transition.

\subsection{Critical-point interchanges and the Frobenius axioms}

A critical-point interchange corresponds to reordering the levels at which two elementary cobordisms are composed. We enumerate the cases.

\begin{lemma}[Interchange moves]\label{lem:interchange-moves}
Let $C_1$ and $C_2$ be two elementary generators whose critical points have adjacent heights. If $C_1$ and $C_2$ act on disjoint boundary circles, then interchanging their heights amounts to rewriting $(C_1 \sqcup \id) \circ (\id \sqcup C_2)$ as $(\id \sqcup C_2) \circ (C_1 \sqcup \id)$. This identity holds in any symmetric monoidal category.

If $C_1$ and $C_2$ share boundary circles, the interchange corresponds to one of the following identities:
\begin{enumerate}
\item Two multiplications sharing an output: associativity \eqref{eq:associativity}.
\item Two comultiplications sharing an input: coassociativity \eqref{eq:coassociativity}.
\item A multiplication followed by a comultiplication on the same circle: the Frobenius relation \eqref{eq:frobenius-relation}.
\item A cap followed by a multiplication: the unit relation \eqref{eq:unitality-left}.
\item A comultiplication followed by a cup: counitality \eqref{eq:counitality}.
\item Swapping the two inputs of a multiplication: commutativity \eqref{eq:commutativity}.
\end{enumerate}
\end{lemma}

\begin{proof}
Each case is verified by identifying the two Morse presentations with the left- and right-hand sides of the corresponding algebraic relation. Cases (1)--(2) correspond to two ways of decomposing a four-holed sphere. Case (3) is the two decompositions of the four-holed sphere treated in Proposition~\ref{prop:cob2-relations}(5). Cases (4)--(5) are the cancellation of a cap or cup against a pair of pants. Case (6) uses the orientation-preserving diffeomorphism that exchanges the two incoming circles of the pair of pants.
\end{proof}

\subsection{Birth--death transitions and unit/counit cancellation}

A birth--death transition creates or annihilates a cancelling pair of critical points. In two dimensions, the relevant pairs are:
\begin{itemize}
\item Index 0 and index 1: a cap composed with one leg of a pair of pants.
\item Index 1 and index 2: one leg of a reversed pair of pants composed with a cup.
\end{itemize}

\begin{lemma}[Birth--death invariance]\label{lem:birth-death}
Let a birth transition introduce a cap followed immediately by a merge (index-0/index-1 pair). The resulting cobordism is diffeomorphic to the identity cylinder. Under $Z_V$, this becomes
\begin{equation}\label{eq:birth-death-0-1}
m \circ (\eta \otimes \id_V) = \id_V,
\end{equation}
which is the left-unitality relation \eqref{eq:unitality-left}. Similarly, an index-1/index-2 birth gives
\begin{equation}\label{eq:birth-death-1-2}
(\epsilon \otimes \id_V) \circ \Delta = \id_V,
\end{equation}
the left-counitality relation \eqref{eq:counitality}. The death transitions yield the same identities read in reverse.
\end{lemma}

\begin{proof}
A cancelling pair of adjacent critical points produces a cobordism that is diffeomorphic to the product $S^1 \times [0,1]$ relative to the boundary. Functoriality then requires $Z_V$ to assign the identity, which is exactly the content of the unit and counit relations.
\end{proof}

\subsection{Conclusion of the well-definedness proof}

\begin{theorem}[Well-definedness of $Z_V$]\label{thm:ZV-well-defined}
Let $(V,m,\eta,\epsilon)$ be a commutative Frobenius algebra. The assignment $Z_V$ defined on generators by \eqref{eq:ZV-circle}--\eqref{eq:ZV-pants} extends to a well-defined symmetric monoidal functor
\begin{equation}\label{eq:ZV-functor}
Z_V : \mathrm{Cob}_2 \longrightarrow \mathrm{Vect}_\C.
\end{equation}
In particular, for any cobordism $M$, the linear map $Z_V(M)$ is independent of the Morse function used to decompose $M$.
\end{theorem}

\begin{proof}
Let $f_0$ and $f_1$ be two Morse functions on $M$ with distinct critical values. By Theorem~\ref{thm:cerf}, they are connected by a generic path $\{f_t\}$ passing through finitely many transitions, each of which is either a critical-point interchange or a birth--death event.

By Lemma~\ref{lem:interchange-moves}, invariance under critical-point interchanges follows from commutativity, associativity, coassociativity, unitality, counitality, and the Frobenius relation---precisely the axioms of a commutative Frobenius algebra.

By Lemma~\ref{lem:birth-death}, invariance under birth--death transitions follows from unitality and counitality.

Since every transition preserves the value of $Z_V$, the endpoints $f_0$ and $f_1$ yield the same linear map. Because $f_0$ and $f_1$ were arbitrary, $Z_V(M)$ depends only on the diffeomorphism class of $M$ relative to the boundary, and $Z_V$ is a well-defined functor.
\end{proof}

\begin{remark}[Relation to the generators-and-relations presentation]\label{rem:generators-relations}
The proof above is equivalent to showing that the Frobenius relations form a complete set of relations for $\mathrm{Cob}_2$. This is the presentation theorem: $\mathrm{Cob}_2$ is the free symmetric monoidal category on a single commutative Frobenius algebra object \cite[Theorem~3.6.19]{Kock2004Frobenius}. Cerf theory provides the bridge between the topological statement (any two Morse decompositions are related by moves) and the algebraic statement (the moves are exactly the Frobenius axioms).
\end{remark}


\section{Mapping-class invariance}\label{appsec:mapping-class}

\subsection{Statement of the invariance property}

For a closed oriented surface $\Sigma_g$, viewed as a cobordism $\emptyset \to \emptyset$, the TQFT assigns a number $Z_V(\Sigma_g) \in \C$. The well-definedness theorem above guarantees that this number is a diffeomorphism invariant. In particular:

\begin{proposition}[Dependence on genus alone]\label{prop:genus-only}
Let $(V,m,\eta,\epsilon)$ be a commutative Frobenius algebra. The partition function $Z_V(\Sigma_g)$ depends only on the genus $g$, not on any additional marking or decomposition of the surface.
\end{proposition}

\begin{proof}
The classification of closed oriented surfaces states that every such surface is diffeomorphic to exactly one $\Sigma_g$. Theorem~\ref{thm:ZV-well-defined} guarantees that $Z_V$ assigns the same value to diffeomorphic surfaces. Hence $Z_V(\Sigma_g)$ is determined by the genus.
\end{proof}

\subsection{Dehn twists act trivially}

A stronger version of the invariance statement concerns surfaces with boundary. For a surface $\Sigma_{g,n}$ of genus $g$ with $n$ boundary circles, the mapping class group $\mathrm{Mod}(\Sigma_{g,n})$ acts on the vector space $Z_V(S^1)^{\otimes n}$ associated to the boundary. In a general TQFT, this action can be nontrivial. In a 2D TQFT built from a commutative Frobenius algebra, however, the action is always trivial.

\begin{theorem}[Triviality of the mapping-class-group action]\label{thm:mcg-trivial}
Let $Z_V$ be the 2D TQFT associated to a commutative Frobenius algebra $(V,m,\eta,\epsilon)$. For any compact oriented surface $\Sigma$ with boundary, the mapping class group $\mathrm{Mod}(\Sigma)$ acts trivially on $Z_V(\Sigma)$.
\end{theorem}

\begin{proof}
The mapping class group of a surface is generated by Dehn twists \cite{Dehn1938,Lickorish1962}. It suffices to show that every Dehn twist acts as the identity.

\medskip

\noindent\emph{Step 1: Dehn twists along non-separating curves.}
Let $\gamma \subset \Sigma$ be a simple closed non-separating curve. The effect of a Dehn twist $T_\gamma$ on the cobordism structure is seen by cutting $\Sigma$ along $\gamma$. This introduces a pair of boundary circles on which one glues back with a twist. The local picture near the cut is a cylinder $S^1 \times [0,1]$ in which the two ends are identified after applying the twist diffeomorphism.

Under the TQFT functor, the cylinder maps to $\id_V$ (the identity on $V$). The Dehn twist replaces this cylinder with a ``twisted cylinder.'' But a Dehn twist along the core of a cylinder is an orientation-preserving diffeomorphism that is the identity on the boundary. The Alexander trick shows that any such diffeomorphism of the annulus that fixes the boundary pointwise is isotopic to the identity relative to the boundary. Hence the induced map on $V$ is the identity.

More precisely, when the curve $\gamma$ is cut open, the local contribution to $Z_V$ involves multiplying and then comultiplying along the same circle. The result is:
\begin{equation}\label{eq:dehn-twist-local}
Z_V(T_\gamma) = m \circ \tau \circ \Delta = m \circ \Delta = H,
\end{equation}
where the second equality uses commutativity of $m$ (equation \eqref{eq:commutativity}). But the \emph{untwisted} cobordism gives the same map $m \circ \Delta = H$. Since commutativity ensures $m \circ \tau = m$, the twist has no effect.

\medskip

\noindent\emph{Step 2: Dehn twists along separating curves.}
If $\gamma$ separates $\Sigma$ into two pieces $\Sigma_1$ and $\Sigma_2$ with $\gamma$ as common boundary, then
\begin{equation}\label{eq:separating-dehn}
Z_V(\Sigma) = Z_V(\Sigma_2) \circ Z_V(\Sigma_1).
\end{equation}
A Dehn twist along $\gamma$ inserts a twisted cylinder between $\Sigma_1$ and $\Sigma_2$. By the same argument as above, the twisted cylinder acts as the identity on $V$, so the Dehn twist acts trivially.

\medskip

\noindent Since both non-separating and separating Dehn twists act trivially, and these generate the full mapping class group, the action of $\mathrm{Mod}(\Sigma)$ on $Z_V(\Sigma)$ is trivial.
\end{proof}

\begin{remark}[Why commutativity is essential]\label{rem:commutativity-essential}
The key step in the proof is equation~\eqref{eq:dehn-twist-local}, where commutativity $m \circ \tau = m$ ensures that reordering the two inputs to the multiplication has no effect. For a \emph{non-commutative} Frobenius algebra, the Dehn twist would act nontrivially, and the resulting theory would not be a TQFT in the oriented sense---it would instead be an \emph{open-closed} TQFT or carry additional structure (e.g., an area-dependent weighting). This is why the classification theorem requires commutativity: the mapping class group must act trivially for the partition function to depend only on diffeomorphism type.
\end{remark}

\begin{remark}[The partition function revisited]\label{rem:partition-revisited}
Combining Theorem~\ref{thm:mcg-trivial} with Theorem~\ref{thm:genus-g-handle}, the partition function of any closed surface is
\begin{equation}\label{eq:partition-genus-only}
Z_V(\Sigma_g) = \epsilon(e^g) = \sum_{i=1}^r \theta_i^{\,1-g}
\end{equation}
in the semisimple case, where the $\theta_i = \epsilon(p_i)$ are the weights of orthogonal idempotents. This formula is manifestly independent of any choices beyond the genus and the Frobenius algebra data.
\end{remark}


\section{Handle-operator spectra for finite groups}\label{appsec:handle-spectrum}

We now specialize to the group-algebra Frobenius algebra $A_G = Z(\C[G])$ of Example~\ref{ex:center-group-algebra} and compute the spectrum of the handle operator explicitly. This yields a closed-form expression for the genus-$g$ partition function in terms of group-theoretic data.

\subsection{General finite-group formula}

Recall from Example~\ref{ex:center-group-algebra} that $A_G$ has basis given by the class sums
\begin{equation}\label{eq:class-sums-app}
z_i = \sum_{g \in \mathcal{C}_i} g, \qquad i = 1, \dots, r,
\end{equation}
and counit $\epsilon_G(\sum_g a_g g) = a_e$. The center $A_G$ is semisimple (since $\C[G]$ is semisimple by Maschke's theorem), so it decomposes into a direct sum of one-dimensional algebras spanned by orthogonal idempotents.

\begin{proposition}[Primitive idempotents of the center of the group algebra]\label{prop:ZG-idempotents}
The primitive orthogonal idempotents of $A_G = Z(\C[G])$ are
\begin{equation}\label{eq:idempotent-formula}
p_\rho = \frac{\dim \rho}{|G|} \sum_{g \in G} \chi_\rho(g^{-1})\, g,
\end{equation}
where $\rho$ ranges over the irreducible representations of $G$ and $\chi_\rho$ is the character of $\rho$.
\end{proposition}

\begin{proof}
This is the standard decomposition of the identity in the center of a semisimple group algebra. The elements $p_\rho$ are the central idempotents corresponding to the isotypic components of the regular representation. They satisfy $p_\rho p_\sigma = \delta_{\rho\sigma} p_\rho$ by orthogonality of characters.
\end{proof}

\begin{proposition}[Handle-operator spectrum for the center of the group algebra]\label{prop:handle-spectrum-G}
With the counit $\epsilon_G$ of \eqref{eq:group-trace}, the handle-operator eigenvalues are
\begin{equation}\label{eq:handle-eigenvalue}
\lambda_\rho = \theta_\rho^{-1} = \frac{|G|}{(\dim \rho)^2},
\end{equation}
where
\begin{equation}\label{eq:theta-rho}
\theta_\rho = \epsilon_G(p_\rho) = \frac{(\dim \rho)^2}{|G|}.
\end{equation}
The handle operator acts on the idempotent basis by
\begin{equation}\label{eq:H-on-p-rho}
H(p_\rho) = \frac{|G|}{(\dim \rho)^2}\, p_\rho.
\end{equation}
\end{proposition}

\begin{proof}
By equation~\eqref{eq:idempotent-formula},
\begin{equation}\label{eq:epsilon-p-rho-calc}
\epsilon_G(p_\rho) = \frac{\dim \rho}{|G|} \chi_\rho(e) = \frac{\dim \rho}{|G|} \cdot \dim \rho = \frac{(\dim \rho)^2}{|G|}.
\end{equation}
Proposition~\ref{prop:idempotent-genus} then gives $H(p_\rho) = \theta_\rho^{-1} p_\rho = \frac{|G|}{(\dim \rho)^2} p_\rho$.
\end{proof}

\begin{theorem}[Genus-$g$ partition function for finite-group gauge theory]\label{thm:genus-g-finite-group}
Let $G$ be a finite group and let $Z_G$ denote the 2D TQFT associated to $(A_G, m, \eta, \epsilon_G)$. Then
\begin{equation}\label{eq:genus-g-group}
Z_G(\Sigma_g) = \sum_{\rho \in \mathrm{Irr}(G)} \left(\frac{(\dim \rho)^2}{|G|}\right)^{1-g} = \sum_{\rho \in \mathrm{Irr}(G)} \left(\frac{|G|}{(\dim \rho)^2}\right)^{g-1}.
\end{equation}
Equivalently, with the finite-gauge-theory normalization $\widetilde{\epsilon}_G = \frac{1}{|G|}\epsilon_G$,
\begin{equation}\label{eq:genus-g-gauge-normalized}
\widetilde{Z}_G(\Sigma_g) = \frac{|\mathrm{Hom}(\pi_1(\Sigma_g), G)|}{|G|} = |G|^{2g-2} \sum_{\rho \in \mathrm{Irr}(G)} \frac{1}{(\dim \rho)^{2g-2}}.
\end{equation}
\end{theorem}

\begin{proof}
The first line follows directly from Proposition~\ref{prop:idempotent-genus} applied with $\theta_\rho = (\dim\rho)^2/|G|$. For the second formula, the rescaling $\widetilde{\epsilon}_G = \frac{1}{|G|}\epsilon_G$ gives $\widetilde{\theta}_\rho = (\dim\rho)^2/|G|^2$, and equation~\eqref{eq:genus-semismple} yields
\begin{equation}\label{eq:normalized-sum}
\widetilde{Z}_G(\Sigma_g) = \sum_{\rho} \left(\frac{(\dim \rho)^2}{|G|^2}\right)^{1-g} = |G|^{2g-2} \sum_{\rho} \frac{1}{(\dim\rho)^{2g-2}}.
\end{equation}

To identify this with the homomorphism count, recall the classical Mednykh formula \cite{Mednykh1978}:
\begin{equation}\label{eq:mednykh}
|\mathrm{Hom}(\pi_1(\Sigma_g), G)| = |G|^{2g-1} \sum_{\rho \in \mathrm{Irr}(G)} \frac{1}{(\dim \rho)^{2g-2}}.
\end{equation}
Dividing by $|G|$ gives \eqref{eq:genus-g-gauge-normalized}.
\end{proof}

\begin{remark}[Interpretation]\label{rem:mednykh-interpretation}
Equation~\eqref{eq:genus-g-gauge-normalized} states that the partition function of untwisted 2D finite gauge theory with gauge group $G$ counts flat $G$-bundles on $\Sigma_g$ (i.e., conjugacy classes of representations $\pi_1(\Sigma_g) \to G$) with each bundle weighted by $1/|\mathrm{Aut}|$. The algebraic origin of this combinatorial identity is the semisimple decomposition of the center of the group algebra.
\end{remark}

\subsection{Explicit computation: \texorpdfstring{$G = \Z_N$}{G = Z\_N}}

For abelian groups, every irreducible representation is one-dimensional, so $\dim \rho = 1$ for all $\rho \in \mathrm{Irr}(\Z_N)$. There are $N$ such representations.

\begin{proposition}[Handle spectrum for $\Z_N$]\label{prop:handle-ZN}
The Frobenius algebra $A_{\Z_N} = \C[\Z_N]$ has $N$ primitive idempotents
\begin{equation}\label{eq:ZN-idempotents}
p_k = \frac{1}{N} \sum_{j=0}^{N-1} \omega^{-jk} g^j, \qquad k = 0, 1, \dots, N-1,
\end{equation}
where $\omega = e^{2\pi i/N}$ and $g$ is the generator of $\Z_N$. Each idempotent has
\begin{equation}\label{eq:ZN-theta}
\theta_k = \epsilon(p_k) = \frac{1}{N},
\end{equation}
and the handle operator has a single eigenvalue $\lambda = N$ with multiplicity $N$:
\begin{equation}\label{eq:ZN-handle-spectrum}
H = N \cdot \id_V.
\end{equation}
The genus-$g$ partition function is
\begin{equation}\label{eq:ZN-genus-g}
Z_{\Z_N}(\Sigma_g) = N^g.
\end{equation}
With the gauge-theory normalization $\widetilde{\epsilon} = \frac{1}{N}\epsilon$,
\begin{equation}\label{eq:ZN-gauge-normalized}
\widetilde{Z}_{\Z_N}(\Sigma_g) = N^{2g-1} = \frac{|\mathrm{Hom}(\pi_1(\Sigma_g), \Z_N)|}{N}.
\end{equation}
\end{proposition}

\begin{proof}
Since $\Z_N$ is abelian, $A_{\Z_N} = Z(\C[\Z_N]) = \C[\Z_N]$. The irreducible representations are the characters $\chi_k(g^j) = \omega^{jk}$, each of dimension 1. Formula~\eqref{eq:idempotent-formula} gives \eqref{eq:ZN-idempotents}. Computing:
\begin{equation}\label{eq:ZN-epsilon-pk}
\epsilon(p_k) = \frac{1}{N}\chi_k(e) = \frac{1}{N}.
\end{equation}
Since all eigenvalues $\lambda_k = \theta_k^{-1} = N$ are equal, $H = N \cdot \id_V$.

For the partition function:
\begin{equation}\label{eq:ZN-partition-calc}
Z_{\Z_N}(\Sigma_g) = \sum_{k=0}^{N-1} \theta_k^{1-g} = \sum_{k=0}^{N-1} N^{g-1} = N^g.
\end{equation}
We verify directly: $\pi_1(\Sigma_g)$ has $2g$ generators with one relation $\prod_{i=1}^g [a_i, b_i] = 1$. For an abelian target, all commutators vanish and the relation is automatic, giving $|\mathrm{Hom}(\pi_1(\Sigma_g), \Z_N)| = N^{2g}$. Dividing by $|\Z_N| = N$ yields $N^{2g-1}$, confirming \eqref{eq:ZN-gauge-normalized}.
\end{proof}

\begin{remark}[Connection to the toric code]\label{rem:toric-code-connection}
The eigenvalue $\lambda = N$ (or $N^2$ in the gauge-normalized version) is the ``quantum dimension'' contribution per handle. For $G = \Z_N$, the ground-state degeneracy of the associated topological phase on $\Sigma_g$ is $N^{2g}$, which equals $|\mathrm{Hom}(\pi_1(\Sigma_g), \Z_N)|$. This is the topological ground-state degeneracy of the $\Z_N$ toric code discussed in Chapter~\ref{sec:topological-order}.
\end{remark}

\subsection{Explicit computation: \texorpdfstring{$G = S_3$}{G = S3}}

The symmetric group $S_3$ is the smallest non-abelian group. It has order $|S_3| = 6$, three conjugacy classes, and three irreducible representations.

\begin{proposition}[Representation theory of $S_3$]\label{prop:S3-reps}
The conjugacy classes and irreducible representations of $S_3$ are:
\begin{center}
\renewcommand{\arraystretch}{1.3}
\begin{tabular}{c c c c}
\hline
Class $\mathcal{C}$ & Size $|\mathcal{C}|$ & Centralizer $|C_G(g)|$ & Representative \\
\hline
$\mathcal{C}_1 = \{e\}$ & 1 & 6 & $e$ \\
$\mathcal{C}_2 = \{(12),(13),(23)\}$ & 3 & 2 & $(12)$ \\
$\mathcal{C}_3 = \{(123),(132)\}$ & 2 & 3 & $(123)$ \\
\hline
\end{tabular}
\end{center}
\medskip
\begin{center}
\renewcommand{\arraystretch}{1.3}
\begin{tabular}{c c c c c}
\hline
Rep $\rho$ & $\dim \rho$ & $\chi_\rho(e)$ & $\chi_\rho((12))$ & $\chi_\rho((123))$ \\
\hline
$\mathbf{1}$ (trivial) & 1 & 1 & 1 & 1 \\
$\mathrm{sgn}$ (sign) & 1 & 1 & $-1$ & 1 \\
$V_{\mathrm{std}}$ (standard) & 2 & 2 & 0 & $-1$ \\
\hline
\end{tabular}
\end{center}
\end{proposition}

\begin{proposition}[Handle spectrum for $S_3$]\label{prop:handle-S3}
The Frobenius algebra $A_{S_3} = Z(\C[S_3])$ is three-dimensional with primitive idempotents:
\begin{align}
p_{\mathbf{1}} &= \frac{1}{6}(e + (12) + (13) + (23) + (123) + (132)), \label{eq:S3-p-trivial} \\[1mm]
p_{\mathrm{sgn}} &= \frac{1}{6}(e - (12) - (13) - (23) + (123) + (132)), \label{eq:S3-p-sign} \\[1mm]
p_{\mathrm{std}} &= \frac{2}{6}(2e - (123) - (132)). \label{eq:S3-p-std}
\end{align}
The handle-operator eigenvalues are:
\begin{equation}\label{eq:S3-eigenvalues}
\lambda_{\mathbf{1}} = \frac{6}{1} = 6, \qquad
\lambda_{\mathrm{sgn}} = \frac{6}{1} = 6, \qquad
\lambda_{\mathrm{std}} = \frac{6}{4} = \frac{3}{2}.
\end{equation}
\end{proposition}

\begin{proof}
Formula~\eqref{eq:idempotent-formula} gives the idempotents directly from the character table. The eigenvalues follow from \eqref{eq:handle-eigenvalue}: $\lambda_\rho = |G|/(\dim \rho)^2$.
\end{proof}

\begin{proposition}[Genus-$g$ partition function for $S_3$]\label{prop:S3-genus-g}
With the counit $\epsilon_{S_3}$,
\begin{equation}\label{eq:S3-partition}
Z_{S_3}(\Sigma_g) = \frac{1}{6^{g-1}}\left(1 + 1 + 4^g \cdot 6^{1-g}\right) = 6^{1-g} + 6^{1-g} + \left(\frac{3}{2}\right)^{g-1}.
\end{equation}
More explicitly:
\begin{equation}\label{eq:S3-partition-explicit}
Z_{S_3}(\Sigma_g) = 2 \cdot 6^{1-g} + \left(\frac{3}{2}\right)^{g-1}.
\end{equation}
With the gauge-theory normalization $\widetilde{\epsilon} = \frac{1}{6}\epsilon$,
\begin{equation}\label{eq:S3-gauge-partition}
\widetilde{Z}_{S_3}(\Sigma_g) = \frac{|\mathrm{Hom}(\pi_1(\Sigma_g), S_3)|}{6} = 6^{2g-2}\left(\frac{2}{1} + \frac{1}{2^{2g-2}}\right).
\end{equation}
\end{proposition}

\begin{proof}
Applying equation~\eqref{eq:genus-g-group}:
\begin{align}
Z_{S_3}(\Sigma_g) &= \sum_{\rho} \theta_\rho^{1-g} = \left(\frac{1}{6}\right)^{1-g} + \left(\frac{1}{6}\right)^{1-g} + \left(\frac{4}{6}\right)^{1-g} \label{eq:S3-calc-start} \\
&= 6^{g-1} + 6^{g-1} + \left(\frac{3}{2}\right)^{g-1} \nonumber \\
&= 2 \cdot 6^{g-1} + \left(\frac{3}{2}\right)^{g-1}. \label{eq:S3-calc-end}
\end{align}
For the gauge-theory normalization, use $\widetilde{\theta}_\rho = (\dim\rho)^2/36$:
\begin{equation}\label{eq:S3-gauge-calc}
\widetilde{Z}_{S_3}(\Sigma_g) = 2 \cdot 36^{g-1} + \left(\frac{36}{4}\right)^{g-1} = 2 \cdot 6^{2g-2} + \frac{6^{2g-2}}{2^{2g-2}}.
\end{equation}
\end{proof}

\begin{example}[Low-genus checks for $S_3$]\label{ex:S3-low-genus}
\begin{enumerate}
\item \emph{Torus ($g=1$):} $Z_{S_3}(T^2) = 2 + 1 = 3$, which equals $\dim A_{S_3} = r = 3$ (the number of conjugacy classes), consistent with Proposition~\ref{prop:torus-partition}.

With gauge normalization: $\widetilde{Z}_{S_3}(T^2) = \frac{|\{(a,b) \in S_3^2 : ab = ba\}|}{6}$. The number of commuting pairs in $S_3$ is $\sum_{g \in S_3} |C_{S_3}(g)| = 1\cdot 6 + 3\cdot 2 + 2\cdot 3 = 18$, so $\widetilde{Z}_{S_3}(T^2) = 18/6 = 3$. Consistent.

\item \emph{Genus 2 ($g=2$):}
\begin{equation}\label{eq:S3-genus-2}
Z_{S_3}(\Sigma_2) = 2 \cdot 6 + \frac{3}{2} = \frac{27}{2}.
\end{equation}
With gauge normalization: $\widetilde{Z}_{S_3}(\Sigma_2) = \frac{|\mathrm{Hom}(\pi_1(\Sigma_2), S_3)|}{6}$. The Mednykh formula gives
\begin{equation}\label{eq:S3-genus-2-mednykh}
|\mathrm{Hom}(\pi_1(\Sigma_2), S_3)| = 6^3\left(\frac{1}{1} + \frac{1}{1} + \frac{1}{4}\right) = 216 \cdot \frac{9}{4} = 486.
\end{equation}
So $\widetilde{Z}_{S_3}(\Sigma_2) = 486/6 = 81$. We verify: $\widetilde{Z}_{S_3}(\Sigma_2) = 2 \cdot 36 + 9 = 81$. Consistent.

\item \emph{Sphere ($g=0$):} $Z_{S_3}(S^2) = 2 \cdot 6^{-1} + (3/2)^{-1} = 1/3 + 2/3 = 1$. This equals $\epsilon(1_V)$: since $1_V = e \in S_3$ and $\epsilon_G(e) = 1$ (the coefficient of the identity in $e$ is 1). Consistent with Proposition~\ref{prop:sphere-from-disk-gluing}.
\end{enumerate}
\end{example}

\subsection{General pattern and asymptotic behavior}

\begin{remark}[Large-genus asymptotics]\label{rem:large-genus}
For any finite group $G$, the genus-$g$ partition function is dominated at large $g$ by the term with the largest eigenvalue $\lambda_\rho = |G|/(\dim\rho)^2$. This maximum occurs for the trivial representation (where $\dim\rho = 1$), giving $\lambda_{\mathbf{1}} = |G|$. If $G$ is non-abelian, the trivial and sign representations (or more generally all one-dimensional representations) contribute equally large eigenvalues, while higher-dimensional representations are suppressed. Thus
\begin{equation}\label{eq:large-genus-asymptotics}
Z_G(\Sigma_g) \sim |\mathrm{Irr}_1(G)| \cdot |G|^{g-1} \qquad \text{as } g \to \infty,
\end{equation}
where $|\mathrm{Irr}_1(G)|$ is the number of one-dimensional representations of $G$, which equals $|G/[G,G]|$, the order of the abelianization.
\end{remark}

\begin{remark}[Class-sum basis and handle action]\label{rem:class-sum-handle}
One can also express the handle operator in the class-sum basis $\{z_i\}$ rather than the idempotent basis. The result is
\begin{equation}\label{eq:handle-class-sum}
H(z_i) = \frac{|G|}{|\mathcal{C}_i|}\, z_i \cdot (\text{coefficient depending on structure constants}),
\end{equation}
but this basis does not diagonalize $H$ in general. The idempotent basis is the natural eigenbasis because it diagonalizes all central elements simultaneously.
\end{remark}
